% English version of paper/ja_v4/06_限界.md (ver4, 2026-11-17): five paragraphs; the answer-bias paragraph moved to Supplementary Information S2. Wording carried over from v3_08_limitations.tex.
\section{Limitations}\label{sec:limitations}

This study has five limitations.

\textbf{Settings, corpora and encoders.} Section~\ref{sec:res-remaining} uses the all-source setting, and Sects.~\ref{sec:res-methods} and \ref{sec:res-margin}--\ref{sec:res-retrain} the LOCO setting. IEMOCAP in the LOCO setting is only its improvised part, and all corpora are in English. One frozen encoder per modality was used. Fine-tuned audio encoders have been reported to give better dimensional emotion prediction (CCC, the concordance correlation coefficient) than frozen features \cite{wagner2023}, so whether the gains of distillation or the reduction of the error are the same with fine-tuning is unknown. For visual features, the median approximation error of the direction loss was 0.47 times the median direction loss, failing the pre-set criterion (below 1/3; Supplementary Information S2), so the visual rows of Table~\ref{tab:remaining} should be read as values that include errors of the same size as the approximation. The values should be read as values for four or five English corpora and linear classifiers on one frozen encoder per modality. Checking other corpora and fine-tuned encoders is the next step.

\textbf{Analyses and methods built on the same data.} The analysis with other pseudo-label providers (Sect.~\ref{sec:res-margin}) and the comparison with the strongest single classifier (Sect.~\ref{sec:res-prop}) were set up after the computation from saved outputs, without prior predictions. The forms, grids and selection procedures of distillation and the proposal were built and inspected on the evaluation corpora as described in Sect.~\ref{sec:protocol}, and the test folds of each corpus also served as provisional-target tests for other targets. Differences from the product and other combinations may therefore be biased in favour of the proposal. Evaluating once on a new corpus with the procedure fixed is the next step.

\textbf{Retrained classifiers and the classifier's own modality in the pseudo-labels.} Only AVT (AV for CREMA-D) was retrained on the proposal's pseudo-labels, not the audio classifier, for which distillation gave positive gains on all four corpora. The pseudo-labels of distillation also include the retrained classifier's own modality, so the gains cannot be attributed to the other modalities alone. The effect of the proposal should therefore be read as values for the fused predictions and retrained AVT. Retraining the audio, visual and text classifiers on the proposal's pseudo-labels, and on pseudo-labels without the classifier's own modality, is the next step.

\textbf{Shared content.} Speakers are separated, but the same person (MOSEI), episode (MELD) or storyline (IEMOCAP) may appear on both sides, and in CREMA-D the same 12 sentences appear on both sides. How much agreement from shared content enters UAR is therefore unknown, and CREMA-D values in particular should be read as values with the same sentences on both sides. Repeating the comparisons with splits that also separate sentences and episodes is the next step.

\textbf{Estimating the advantage, and intervals.} The advantage was measured with all candidate labels, which a new corpus lacks, and only at K = 5 and 25. Section~\ref{sec:res-margin} therefore describes the variants whose UAR rose rather than giving a procedure for choosing pseudo-label providers. Intervals are not corrected for multiple comparisons, and counts of conditions are not tests. Estimating the advantage within the K labels and comparing it with the value measured on all candidate labels is the next step.
